% !TeX encoding = UTF-8 !TeX spellcheck = en_US
\documentclass[11pt, a4paper]{article}

% The title of the current document to be produced.
\newcommand{\doctitle}{Assignment \#6}
\newcommand{\assignments}{\bluetext{\textbf{Assignment:}}}
%
\setlength{\unitlength}{1in}
\renewcommand{\arraystretch}{2}
\setlength{\columnsep}{2cm}

\newcommand\dunderline[3][-1pt]{{%
  \sbox0{#3}%
  \ooalign{\copy0\cr\rule[\dimexpr#1-#2\relax]{\wd0}{#2}}}}

\usepackage[sfdefault]{atkinson}
\usepackage[T1]{fontenc}

\def\code#1{\bluetext{\texttt{#1}}}


%------------------------------------------------------------
% Import commands for both teacher and course information.  | NOTE: Change your
% teacher and course info in these files. |
%------>------>------>------>------>------>------>------>-->|
%
%------------------------------------------------------------
%-- Import packages and custom command definitons.          |
%------>------>------>------>------>------>------>------>-->|
%% ==================================
%% ===== Teacher info             ===
%% ==================================
%%
\newcommand{\instructor}{Katelyn Breivik}
\newcommand{\office}{Wean 8319}
\newcommand{\hours}{3pm on Mondays}
\newcommand{\phone}{}
\newcommand{\college}{}
\newcommand{\email}{kbreivik@andrew.cmu.edu}
\newcommand{\faculty}{}
\newcommand{\department}{Department of Physics}
                              %|
%% ==================================
%% ===== Course-specific commands ===
%% ==================================

%- Instructions: change course info here. 
\newcommand{\semester}{Fall 2023}

\newcommand{\ponderation}{2-4-3 (Theory-Lab-Homework)}
\newcommand{\coursetitle}{Astrophysics of Stars and the Galaxy}
\newcommand{\coursenumber}{[33-467]}
\newcommand{\prerequisite}{33-224, 33-234}
 
\input{../includes/packages-imports}                          %|  
\input{../includes/custom-commands}   
%
%---> Genereate & inject metadata                           |
\input{../includes/hyperef.doc.info}                          %|
%------------------------------------------------------------

\topmargin -70pt
\begin{document} 

\normalfont

%-------------------------------------------------------------
%-- Make the header of the document                          |
%------>------>------>------>------>------>------>------>--> |
%---------------------------------------------------------------------
%-- The following produces the document header including the title.  |
%---------------------------------------------------------------------

\noindent % <-- need to have this first.

\hfill	
\begin{minipage}{0.60\textwidth}%
	\raggedleft%
	{\Large \textsc{\doctitle}\par}
	\doublerule 
\end{minipage}%
\vspace{0.9cm}

  
%-------------------------------------------------------------
%-- Problem # 1                                              |
%------>------>------>------>------>------>------>------>--> |
\noindent \dunderline{1.0pt}{\textbf{Problem \#1: Binary Evolution (36
points)}}\\
We have officially covered all the top level things necessary to understand the
general principles of binary evolution! This problem asks you to prove it.
\vspace{8pt}

\noindent \textbf{Problem \#1a: Install and run COSMIC (12 pts)}\\
\noindent COSMIC is a binary population synthesis code that can evolve binary
stars from ZAMS up to the formation and merger of the stellar remnants in the
binary. The first part of this problem is to install and import COSMIC into a
jupyter notebook. Since COSMIC only builds for linux and mac os, if you have a
windows machine, you should use Google Colab (which runs linux). You can follow
the installation instructions in the
\href{https://cosmic-popsynth.github.io/COSMIC/install/index.html}{COSMIC
documentation}
\vspace{4pt}

\noindent Once you have COSMIC installed, run the
\href{https://cosmic-popsynth.github.io/COSMIC/examples/index.html}{first
example} in the documentation which evolves a massive binary star that produces
black holes which kind of match the mass of GW150914. To prove that you have run
the code, print the 'tphys', 'mass\_1', 'mass\_2', 'kstar\_1', 'kstar\_2',
'porb', 'ecc', and 'evol\_type' columns from the bpp DataFrame. You can use the
`utils.convert\_kstar\_evol\_type' function in COSMIC by first importing utils,
then calling the function with your bpp DataFrame as an argument as shown in the
documentation. You can also read about the kstar and evol\_type keys in the
\href{https://cosmic-popsynth.github.io/COSMIC/output_info/index.html}{COSMIC
documentation}.

\vspace{4pt} 

\noindent Describe the evolution of the binary in words. Is the Roche lobe of
either component in the binary filled? If so, how does the Roche overflow affect
the orbit? What about the masses of each component in the binary? How does the
formation of each black hole alter the orbit?
\vspace{8pt}

\noindent \textbf{Problem \#1b: Low mass binary evolution (12 pts)}\\
\noindent Evolve a ZAMS binary with the following parameters:\newline
$mass\_1=6.0\,\rm{M_{\odot}}$, $mass\_2=1.5\,\rm{M_{\odot}}$,
$porb=600\,\rm{days}$, $ecc=0$, $Z=0.014$\newline
In the `Evolve.evolve' function call, supply the keyword argument: `dtp=0.0'.
This will tell COSMIC to keep all timesteps for the binary in the `bcm'
DataFrame. 

\vspace{4pt}
\noindent Once the binary is evolved, print the bpp DataFrame with the same
columns as in Problem 1a. Next, make a figure of the orbital period vs time and
a second figure of the two component masses vs time. Be sure to label each axis,
and include a legend for the mass figure. If most of the evolution occurs at
certain times, you should trim the x limits of your figure.  

\vspace{4pt}
\noindent Similar to Problem 1a, describe the evolution of the binary in words
taking care to describe the reason for all the different features in each of
your figures. 

\vspace{10pt}
\noindent \dunderline{1.0pt}{\textbf{Problem \#1c: High mass binary evolution
(12 pts)}}\\
\noindent Since natal kicks are implemented as a stochastic (random) process, if
you want to reproduce the exact evolution of a binary, you'll need to supply a
random seed to numpy. For this problem, you should seed numpy's random number
generator with 5. Once the seed is set, evolve a ZAMS binary with the following
parameters:\newline
$mass\_1=50.0\,\rm{M_{\odot}}$, $mass\_2=35.0\,\rm{M_{\odot}}$,
$porb=120\,\rm{days}$, $ecc=0$, $Z=0.014$\newline
In the `Evolve.evolve' function call, supply the keyword argument: `dtp=0.0'. As
before, this will tell COSMIC to keep all timesteps for the binary in the `bcm'
DataFrame. 

\vspace{4pt}
\noindent Once the binary is evolved, print the bpp DataFrame with the same
columns as in Problem 1a. Next, make a figure of the orbital period vs time and
a second figure of the two component masses vs time. Be sure to label each axis,
and include a legend for the mass figure. If most of the evolution occurs at
certain times, you should trim the x limits of your figure.  

\vspace{4pt}
\noindent Similar to Problem 1a, describe the evolution of the binary in words
taking care to describe the reason for all the different features in each of
your figures. What is the main outcome difference from this binary and the
binary in problem 1a?

\vspace{20pt}

\noindent \dunderline{1.0pt}{\textbf{Problem \#2: Gravitational wave analysis (6
points)}}\\
The year is 2038. Prof. Breivik is as hype as she's ever been because the first
data from LISA has been released. One signal has been measured to have a
dimensionless strain value of $h=1.8\times10^{-22}$, a gravitational-wave
frequency of $f_{\rm{GW}}=2.75\,\rm{mHz}$ and a gravitational-wave frequency
evolution of $\dot{f}_{\rm{GW}}=1.1 \times 10^{-16}\rm{Hz/s}$. 

\vspace{4pt}
\noindent Based on these measured quantities, solve for the chirp mass,
$\mathcal{M}_c$, and distance, $D$ to this new source. What type of source do
you expect is producing this gravitational-wave signal based on it's chirp mass?



\end{document} 
